\section{Introduction}

When several layers share a learned bank of parameters, their effective
weights can be written as $\Theta=AB$: each row of a router $A$ selects a
convex combination of the rows of a jointly learned basis matrix $B$.
The router is useful for optimization and compression, but its entries are
not uniquely determined by the effective weights. What additional observation
would identify the \emph{current} factors, rather than merely select a
convenient representation?

Static ambiguity is established simplex-structured matrix factorization
(SSMF) theory~\citep{vuthanh2023bounded}. For example, the invertible
transform $M_s=(1-s)I+s\ones\ones^\top/K$, $0<s<1$, preserves the product
under $(A,B)\mapsto(AM_s,M_s^{-1}B)$ while moving each router row toward
uniformity. This increases the entropy of nonuniform rows. It preserves
coordinate ordering and argmax assignments; other feasible, nonuniform
changes of basis can change the argmax partition
(Appendix~\ref{app:secondary-theory-statements}). An effective model alone
therefore need not determine coordinate-based summaries or decisions.

We study an additional observation supplied by a specified optimizer:
the instantaneous response of the effective parameters to an arbitrary
effective gradient. With $A=\operatorname{softmax}_{\rm row}(Z/\tau)$,
fixed temperature, and known positive Euclidean learning rates for $Z$ and
$B$, this is a linear map $\mathcal R_{A,B}:G\mapsto-\dot\Theta$.
It describes controlled interventions at a fixed checkpoint. The full map contains more information than a single task-gradient
response at that checkpoint. It is a controlled observation distinct from
a passive training trajectory and fixes the optimizer metric being compared.

Our central result is that this observation resolves the static ambiguity.
For strictly positive, full-column-rank routers and full-row-rank bases
with the same product, equal response operators imply a common basis
permutation, and conversely. The proof recovers a Gram matrix and a family
of router covariances from the response, then reduces the remaining
ambiguity to simultaneous diagonalization. On uniformly nondegenerate
domains, the same chain gives an explicit inverse-stability bound.

The complete-operator premise can also be replaced by finite, imperfect
observations. If $D-K\ge L$, we design $K$ unit-norm queries from a noisy
observed product. Every pair of factors consistent with the product and
response error budgets is close modulo permutation, with an explicit bound
on the diameter of the \emph{entire} feasible set. The true basis need not
lie in the estimated product subspace, and candidates need not belong to one
local solution branch. This is an information guarantee on a declared domain,
not an efficient recovery algorithm or a data-derived certificate of that
domain's rank and interior margins.

The contributions are:
\begin{enumerate}
\item A model-specific characterization of the complete Euclidean-response
stabilizer on a common product fiber, with rank and observation-boundary
counterexamples (Section~\ref{sec:response_identifiability}).
\item An explicit inverse bound for the original factors, a bridge allowing
small product perturbations, and a whole-feasible-set diameter bound from
noisy products and $K$ designed response queries
(Sections~\ref{sec:quantitative}--\ref{sec:finite-observations}).
\item A frozen-protocol audit of shared modules on nine pretrained language
models (160M--8B), across three adapter seeds and three gradient domains.
Recovery, precision failures, transferred rejection rules, and a matched
structural decision test distinguish coordinate identification from
downstream usefulness (Section~\ref{sec:experiments}).
\end{enumerate}

The response, diagonalization, robust-inverse, and probing tools have
substantial precedents. Our contribution is their complete reduction for
this observation model and the resulting finite-noisy guarantee, rather
than a new general symmetry principle. The identified object is the current
optimizer-dependent factorization. Neither the theorem nor the experiments
recover the historical process that produced a model.
